% Standalone problem document, generated from the adjacent README.md.
% Compile with XeLaTeX. Mathematical content is maintained in Markdown.
\documentclass[11pt,a4paper]{article}
\usepackage[fontsize=10.5pt]{scrextend}
\usepackage[margin=22mm,headheight=15pt,headsep=9mm,footskip=11mm]{geometry}
\usepackage{fontspec}
\setmainfont{texgyrepagella}[Extension=.otf,UprightFont=*-regular,BoldFont=*-bold,ItalicFont=*-italic,BoldItalicFont=*-bolditalic]
\setsansfont{texgyreheros}[Extension=.otf,UprightFont=*-regular,BoldFont=*-bold,ItalicFont=*-italic,BoldItalicFont=*-bolditalic]
\setmonofont{texgyrecursor}[Extension=.otf,UprightFont=*-regular,BoldFont=*-bold,ItalicFont=*-italic,BoldItalicFont=*-bolditalic,Scale=MatchLowercase]
\usepackage{amsmath,amssymb}
\usepackage{unicode-math}
\setmathfont{texgyrepagella-math.otf}
\usepackage{xcolor,microtype,enumitem,needspace,fancyhdr}
\usepackage{bookmark}
\definecolor{ink}{HTML}{17344B}
\definecolor{accent}{HTML}{197D80}
\definecolor{muted}{HTML}{596773}
\hypersetup{colorlinks=true,linkcolor=accent,urlcolor=accent,pdftitle={FR-05 - Vanishing
injectivity probability at 4d minus 5 complex phase
measurements},pdfauthor={Open Problems in Numerical Linear Algebra}}
\urlstyle{same}
\setlength{\parindent}{0pt}
\setlength{\parskip}{5pt plus 1pt minus 1pt}
\linespread{1.04}
\setlength{\emergencystretch}{3em}
\widowpenalty=10000
\clubpenalty=10000
\displaywidowpenalty=10000
\allowdisplaybreaks[1]
\setlist{nosep,leftmargin=1.5em}
\providecommand{\tightlist}{\setlength{\itemsep}{0pt}\setlength{\parskip}{0pt}}
\providecommand{\pandocbounded}[1]{#1}
\pagestyle{fancy}
\fancyhf{}
\fancyhead[L]{\sffamily\footnotesize\color{muted}OPEN PROBLEMS IN NUMERICAL LINEAR ALGEBRA}
\fancyhead[R]{\sffamily\bfseries\footnotesize\color{accent}FR-05}
\fancyfoot[L]{\parbox[b]{0.9\textwidth}{\sffamily\fontsize{8}{10}\selectfont\color{muted}Editorial ratings; a bounded search cannot certify that a problem remains unsolved.\\Literature check: 2026-09-08}}
\fancyfoot[R]{\sffamily\footnotesize\color{muted}\thepage}
\renewcommand{\headrulewidth}{0.3pt}
\renewcommand{\headrule}{\hbox to\headwidth{\color{accent}\leaders\hrule height \headrulewidth\hfill}}
\makeatletter
\renewcommand{\section}{\Needspace{5\baselineskip}\@startsection{section}{1}{0pt}{12pt plus 2pt minus 2pt}{4pt}{\sffamily\bfseries\large\color{ink}}}
\renewcommand{\subsection}{\Needspace{4\baselineskip}\@startsection{subsection}{2}{0pt}{9pt plus 1pt minus 1pt}{3pt}{\sffamily\bfseries\normalsize\color{ink}}}
\makeatother
\setcounter{secnumdepth}{0}
\begin{document}
{\sffamily\small\color{muted}Frames and matrix designs\par}
\vspace{3pt}
{\raggedright\hyphenpenalty=10000\exhyphenpenalty=10000\sffamily\bfseries\fontsize{21}{24}\selectfont\color{ink}Vanishing
injectivity probability at 4d minus 5 complex phase measurements\par}
\vspace{7pt}
{\sffamily\small\textbf{Difficulty:} hard\quad\textbf{Importance:} interesting
to the community\par}
\vspace{4pt}
{\color{accent}\hrule height 0.8pt}
\vspace{6pt}
\textbf{Status:} source-stated conjecture, asymptotic part only; no
resolution found in screening on 2026-09-08.

For every integer \(d\geq2\), draw \(A_d\in\mathbb C^{(4d-5)\times d}\)
with independent standard complex Gaussian entries: the real and
imaginary parts are independent \(N(0,1/2)\). Let \(p_d\) be the
probability that \[
|A_dx|=|A_dy|\quad\Longrightarrow\quad
y=e^{\mathrm i\theta}x\ \text{for some }\theta\in\mathbb R
\] holds for all \(x,y\in\mathbb C^d\), where absolute value is
componentwise. Equivalently, \(p_d\) is the probability that the induced
phase-retrieval map on vectors modulo global phase is injective.

Is \[
\lim_{d\to\infty}p_d=0?
\]

This is part (b) of Vinzant's conjecture as restated in Randomstrasse
101. It distinguishes injective exceptional measurement systems from
typical matrices at a row count just below \(4d-4\). Injectivity is the
basic identifiability requirement before conditioning and stable
numerical inversion can be addressed.

\subsection{References}\label{references}

\begin{enumerate}
\def\labelenumi{\arabic{enumi}.}
\tightlist
\item
  A. S. Bandeira et al., \emph{Randomstrasse 101: Open Problems of
  2025}, arXiv:2603.29571 (2026), Conjecture 19(b).
  \href{https://arxiv.org/html/2603.29571v1}{Paper}.
\item
  C. Vinzant, \emph{A small frame and a certificate of its injectivity},
  SAMPTA 2015, arXiv:1502.04656. The explicit 11-vector frame in
  \(\mathbb C^4\) demonstrates why the below-\(4d-4\) regime cannot
  simply be dismissed. \href{https://arxiv.org/abs/1502.04656}{Paper}.
\item
  Z. Li, \emph{On Injectivity of Phase Retrieval}, arXiv:2606.17922
  (2026). The abstract and main result establish the strict inequality
  \(p_d<1\), explicitly identifying this as the first part of Vinzant's
  conjecture. \href{https://arxiv.org/abs/2606.17922}{Paper}.
\end{enumerate}

\subsection{Status check --- 2026-09-08}\label{status-check-2026-09-08}

Searched ``Vinzant conjecture injectivity probability limit'', ``phase
retrieval 4M-5 2026'', and checked the latest June 2026 paper
abstract/version. Its positive probability of noninjectivity does not
supply an asymptotic lower bound tending to one. The already proved
assertion \(p_d<1\) is deliberately excluded from this problem.
\end{document}
